\section{Описание алгоритмов}
\subsection{Линейный поиск}
Линейный поиск последовательно проверяет элементы массива до обнаружения
нужного значения или до конца массива~\cite{linear}. Предварительная
сортировка для этого не требуется.

Обозначим количество элементов через \texttt{n}, а число шагов линейного
поиска через \texttt{L(n)}. Если значение отсутствует, алгоритм проверяет
все элементы:
\begin{equation}
\mathtt{L}(\mathtt{n})=\mathtt{n}.
\label{eq:linear}
\end{equation}
Если все элементы различны, значение присутствует и каждая его позиция
равновероятна, среднее число шагов равно
\begin{equation}
\overline{\mathtt{L}}(\mathtt{n})=
\frac{1}{\mathtt{n}}\sum_{\mathtt{i}=1}^{\mathtt{n}}\mathtt{i}
=\frac{\mathtt{n}+1}{2}.
\label{eq:average}
\end{equation}
Здесь \texttt{i} --- номер позиции, начиная с единицы.

\subsection{Бинарный поиск}
Бинарный поиск работает с отсортированным массивом. Алгоритм проверяет
середину текущего диапазона и оставляет только ту половину, в которой
может находиться значение~\cite{binary}.

В программе границы диапазона хранятся в переменных \texttt{left}
и \texttt{right}. Индекс среднего элемента \texttt{middle} вычисляется так:
\begin{equation}
\mathtt{middle}=\mathtt{left}+
\left\lfloor\frac{\mathtt{right}-\mathtt{left}}{2}\right\rfloor.
\label{eq:middle}
\end{equation}
Индексы начинаются с нуля. При несовпадении проверенный элемент исключается.

Для непустого массива максимальное число шагов данной реализации равно
\begin{equation}
\mathtt{B}_{\max}(\mathtt{n})=
\lfloor\log_2\mathtt{n}\rfloor+1,\qquad \mathtt{n}\geq 1.
\label{eq:binary}
\end{equation}
Число шагов растёт логарифмически. Формула \eqref{eq:binary}
описывает худший случай; при совпадении поиск заканчивается сразу.

